\documentclass[11pt]{article}

\usepackage[T1]{fontenc}
\usepackage[utf8]{inputenc}
\usepackage{mathpazo}
\usepackage{amsmath,amssymb,amsthm}
\usepackage{geometry}
\usepackage{microtype}
\usepackage{setspace}
\usepackage{titlesec}
\usepackage{hyperref}
\usepackage{enumitem}
\usepackage{booktabs}
\usepackage{array}

\geometry{
  paper=a4paper,
  top=1.25in,
  bottom=1.25in,
  left=1.4in,
  right=1.4in
}

\onehalfspacing

\hypersetup{
  colorlinks=true,
  linkcolor=black,
  citecolor=black,
  urlcolor=black
}

\titleformat{\section}{\large\bfseries}{\thesection.}{0.6em}{}
\titleformat{\subsection}{\normalsize\bfseries}{\thesubsection.}{0.5em}{}

\newtheoremstyle{rthm}
  {10pt}{10pt}{\itshape}{0pt}{\bfseries}{.}{0.5em}{}
\theoremstyle{rthm}
\newtheorem{theorem}{Theorem}[section]
\newtheorem{proposition}[theorem]{Proposition}
\newtheorem{corollary}[theorem]{Corollary}
\newtheorem{lemma}[theorem]{Lemma}

\newtheoremstyle{rdef}
  {8pt}{8pt}{\normalfont}{0pt}{\bfseries}{.}{0.5em}{}
\theoremstyle{rdef}
\newtheorem{definition}[theorem]{Definition}

\theoremstyle{remark}
\newtheorem*{remark}{Remark}

\title{\textbf{The Substance and Accident of Reasoning}\\[6pt]
\large Admissibility, Repair, and the Geometry of Verbosity}
\author{Flyxion\\[4pt]
\small Independent Researcher}
\date{June 2026}

\begin{document}

\maketitle
\thispagestyle{empty}

\begin{abstract}
\noindent
The central result of \textit{Shorter but not Worse}~\cite{bounhar2026}
is often presented as an efficiency finding: models trained on a curriculum
containing moderately easy mathematical problems produce substantially
shorter solutions while preserving or improving reasoning performance.
Yet the deeper significance of the result is not computational but geometric.
The paper reveals a form of representational collapse in which an accidental
property of successful reasoning becomes mistaken for an essential one.
This essay analyses the failure through the lens of admissibility theory.
A representation is admissible when it preserves distinctions required for
future action. The hard-only training curriculum induces inadmissibility
by systematically suppressing evidence that correct solutions can be short.
The resulting policy identifies correctness with length---a projection
collapse in which two properties that ought to be orthogonal become
conflated. The easy-problem repair restores admissibility not by adding
information but by expanding the fiber of successful trajectories to
include compact ones. The information-theoretic argument in the paper
explains why verbosity emerges under length-unpenalised training but
misidentifies the goal: entropy reduction is not intelligence. The verbosity
error is one of three instances of the same deeper mistake---the
identification of a proxy metric with the target itself---alongside
the latent-compression error (lower dimensionality as better
representation) and the reasoning-length error (more steps as
deeper thought). All three are projection collapses in which
accidental properties of training distributions are mistaken for
substantial properties of the task. The solution in each case is
not additional optimisation pressure but representational repair:
introducing training evidence that demonstrates the admissibility
of the target without the proxy.
\end{abstract}

\newpage
\setcounter{page}{1}

%=====================================================================
\section{The Finding and Its Conventional Interpretation}
\label{sec:finding}
%=====================================================================

Reinforcement learning pipelines for mathematical reasoning frequently
discard easy problems on the grounds that they provide little optimisation
signal. When all rollouts for a given problem are correct, the group
relative advantage collapses to zero and the problem contributes no
gradient. Training therefore concentrates on problems of intermediate
and high difficulty whose successful solutions are typically long. Over
time, reward becomes correlated with verbosity, and the model learns to
associate correctness with length.

Bounhar et al.\ demonstrate that reintroducing moderately easy problems
reverses this tendency. Models trained under the mixed curriculum produce
solutions that are nearly half the length of baseline solutions, with
no loss in accuracy on competition-level mathematics and improvements
on several benchmarks. The intervention requires no explicit length
penalty. The brevity is emergent.

The conventional interpretation frames this as a training efficiency
result: easy problems act as implicit length regularisers by providing
stable reward signals associated with short outputs. This framing is
accurate as far as it goes. But it does not explain what has actually
gone wrong in the baseline model or why a particular form of data
augmentation specifically repairs it. To answer these questions requires
a different theoretical vocabulary.

%=====================================================================
\section{Representational Collapse}
\label{sec:collapse}
%=====================================================================

A representation is \emph{admissible} when it preserves distinctions
required for future action. It becomes inadmissible when states with
different consequences are projected into the same representational
region. The hard-only curriculum creates precisely such a projection.

Consider two solutions to a given mathematical problem:

\[
s_1 = \text{short correct proof}, \qquad s_2 = \text{long correct proof}.
\]

The downstream consequence of both trajectories is identical: the
problem is solved and reward~1 is received. Yet the training distribution
systematically exposes the model to successful trajectories of the
second type while suppressing examples of the first. The learned
representation therefore begins to identify correctness with length.
Symbolically, the model's implicit mapping satisfies

\[
F(\text{correct}) \approx F(\text{long})
\]

rather than the admissible condition

\[
F(\text{correct}) \perp F(\text{length}).
\]

The distinction between solution quality and solution verbosity has
collapsed.

This is not merely a statistical artifact. It is a loss of navigational
structure. A planner operating on such a representation can no longer
distinguish between properties that affect reachability and properties
that merely accompany successful trajectories. The resulting policy
cannot reliably separate what is necessary from what is incidental.

%=====================================================================
\section{The Projection Operator and Its Distorted Fibers}
\label{sec:projection}
%=====================================================================

The CLIO framework offers a precise geometric language for this
failure. The hard-only curriculum induces a projection operator

\[
\pi : (\text{quality},\, \text{length}) \;\longrightarrow\; \text{quality}
\]

whose empirical fibers become distorted. In an undistorted representation,
the pre-image of good solutions would contain both short and long
trajectories:

\[
\pi^{-1}(\text{good}) = \{\text{short trajectories}\} \cup \{\text{long trajectories}\}.
\]

Under the hard-only curriculum, the empirical support of this fiber
contracts:

\[
\pi^{-1}(\text{good}) \;\approx\; \{\text{long trajectories}\}.
\]

The fiber's support has collapsed. What was a two-dimensional region in
the quality-length product space has degenerated: the policy assigns
negligible probability mass to the short-trajectory region of the fiber,
placing it effectively outside the reach of the learned representation.

The easy-problem repair operates directly on this support collapse. By
introducing successful trajectories that are both correct and compact,
it repopulates the short-trajectory region of $\pi^{-1}(\text{good})$:

\[
\pi^{-1}(\text{good}) \;\longleftarrow\; \pi^{-1}(\text{good})
\cup \{\text{short correct trajectories}\}.
\]

The fiber's support expands and recovers its extent across the
quality-length space. What appeared to be an
intrinsic feature of successful reasoning---length---is revealed to
be merely one coordinate among many. The projection no longer collapses
the distinction between quality and verbosity.

\begin{definition}[Admissible Reasoning Representation]
A policy's internal representation of solution strategies is
\emph{admissible} if the fibers of the correctness projection contain
both short and long correct trajectories whenever both exist for a
given problem class. Admissibility fails when the fiber contracts to
a single strategy type under training pressure.
\end{definition}

The baseline verbosity-trained model is inadmissible in this sense.
The frugal model, after Stage~1 repair, is admissible: it can locate
correct solutions across the full range of trajectory lengths and
select among them based on problem structure.

%=====================================================================
\section{Entropy Reduction Is Not Intelligence}
\label{sec:entropy}
%=====================================================================

The information-theoretic argument in the paper identifies a genuine
mechanism underlying verbosity. By the chain rule of conditional
entropy:

\[
H(Y \mid X, Z_{t+1}) \;\leq\; H(Y \mid X, Z_t),
\]

where $Y$ is the final answer and $Z_t$ is the length-$t$ prefix
produced by the autoregressive policy. Conditioning on a longer prefix
can only decrease or preserve the conditional entropy of the final
answer. This holds regardless of whether the additional tokens carry
semantic content. In the absence of any length penalty, a policy
rewarded for correct final answers has a weak statistical incentive to
extend its output: verbosity becomes a shortcut for entropy reduction.

This observation is correct but incomplete. It identifies why verbosity
emerges. It does not identify what is wrong with it.

Entropy reduction is not itself intelligence.

A language can reduce uncertainty by merging words---collapsing
distinct meanings into a single form so that the distribution over
interpretations narrows. A map can reduce uncertainty by erasing roads
---reducing the number of possible routes so that the next turn becomes
predictable. A bureaucracy can reduce uncertainty by prohibiting
choices---constraining the action space until outcomes are fully
determined.

In each case the system becomes more predictable. In each case it
simultaneously becomes less useful.

The collapse of distinct meanings into a single form makes
disambiguation harder. The erasure of roads reduces navigational
capacity. The prohibition of choices destroys the agency the system
was designed to support.

The relevant quantity is not entropy but admissibility.

The question is not whether uncertainty decreases after a longer
prefix. The question is whether the distinctions necessary for
navigation survive. A model that reduces uncertainty by making all
successful solutions look long has reduced entropy at the cost of
admissibility. It has made itself more predictable while simultaneously
making itself less capable of producing the full range of correct
solutions its task requires.

From this perspective, the model's verbosity resembles a language
undergoing semantic merger. Distinct reasoning trajectories---short
correct proofs and long correct proofs---are collapsed into a single
category whose defining feature is length. The representation becomes
easier to optimise while simultaneously becoming less faithful to the
structure of the task it is meant to solve.

%=====================================================================
\section{Three Instances of the Same Error}
\label{sec:instances}
%=====================================================================

The verbosity failure belongs to a class of machine learning failures
that share a common structure. In each case, a proxy metric that
correlates with the target in the training distribution becomes
identified with the target itself. The identification holds as long
as the training distribution sustains it. When the distribution
changes, or when the model is deployed outside its training regime,
the identification breaks and the failure becomes visible.

Three instances of this error are worth placing alongside each other.

\paragraph{Latent compression as intelligence.}

The latent fundamentalist mistake is to identify a low-dimensional
representation with a better representation~\cite{flyxion2026latent}.
A model trained to compress input into a compact latent space learns
that lower dimensionality correlates with generalisation in its
training regime. The identification holds when the training
distribution is well-covered and the latent structure captures
genuine variation. It breaks when the latent projection discards
variation that matters for downstream tasks---when it collapses
distinctions required for navigation.

The formal structure is:

\[
F(\text{general representation}) \approx F(\text{low-dimensional}),
\]

when the admissible condition is

\[
F(\text{general representation}) \perp F(\text{dimensionality}).
\]

\paragraph{Entropy reduction as intelligence.}

The verbosity error is to identify lower conditional entropy with
better reasoning. A model trained without length penalties learns
that longer prefixes reduce uncertainty about the final answer.
The identification holds when the training distribution consists
primarily of hard problems where length genuinely correlates with
correctness. It breaks when short correct solutions exist but have
been excluded from the training distribution, leaving the model
unable to produce them.

The formal structure is:

\[
F(\text{correct reasoning}) \approx F(\text{low entropy}),
\]

when the admissible condition is

\[
F(\text{correct reasoning}) \perp F(\text{length}).
\]

\paragraph{Long chains of thought as intelligence.}

The reasoning length error is the most culturally embedded of the
three. Contemporary discourse about AI progress frequently equates
extended chains of thought with deeper reasoning. Test-time
compute scaling is discussed as if more computation were always
better computation. A model that produces ten reasoning steps is
presumed to be engaging in more careful reasoning than one that
produces two, even when the two-step solution is correct and the
ten-step solution is padded.

The formal structure is again the same:

\[
F(\text{good reasoning}) \approx F(\text{many steps}),
\]

when the admissible condition is

\[
F(\text{good reasoning}) \perp F(\text{step count}).
\]

\paragraph{The common structure.}

All three errors have the same form: a genuine correlation in the
training distribution is mistaken for an identity. In the
distribution where the model learned, the proxy and the target
covary. Outside that distribution, they diverge.

The admissibility framework names this structure: it is a projection
collapse. A property of trajectories that is accidental in the
philosophical sense---contingent on the distribution, not
constitutive of the target---has been treated as substantial.
The repair in each case is the same: introduce training examples
that demonstrate the admissibility of the target without the proxy.
Short correct solutions for verbosity. High-dimensional correct
representations for latent compression. Two-step correct proofs
for reasoning length.

What makes these failures persist is that the proxies are not
arbitrary. In many distributions, lower entropy does correlate
with better reasoning. In many distributions, lower dimensionality
does correlate with better generalisation. In many distributions,
longer chains of thought do correlate with correct answers. The
proxies are genuine signals. They become dangerous only when the
training distribution is not representative of the full range of
cases where the target behaviour is admissible.

The model is not failing to learn. It is succeeding in learning
from the data it has. The data is the problem.

\medskip

The common algebraic skeleton of all three cases can be stated
as a single proposition.

\begin{proposition}[Proxy Collapse and Repair]
\label{prop:proxy}
Let $T$ be a target property and $Q$ a proxy such that
$\mathrm{Corr}(T, Q) \approx 1$ in the training distribution
$\mathcal{D}$. A \emph{proxy collapse} occurs when the learner
replaces the distributional co-occurrence

\[
\Pr(T \mid Q) \approx 1
\]

with the stronger identification $T \equiv Q$, treating $Q$ as
constitutive of $T$ rather than merely correlated with it.

The collapse is repaired by introducing training examples in
which $T \land \neg Q$ or $Q \land \neg T$ holds---examples that
demonstrate the non-identity of $T$ and $Q$ directly.

Under this schema, the three cases reduce to:

\medskip
\begin{center}
\renewcommand{\arraystretch}{1.3}
\begin{tabular}{lll}
\toprule
\textbf{Error} & \textbf{Collapse} & \textbf{Repair signal} \\
\midrule
Latent compression & $\text{general} \equiv \text{low-dim}$ &
  high-dim correct representations \\
Verbosity & $\text{correct} \equiv \text{long}$ &
  short correct solutions \\
Reasoning length & $\text{good} \equiv \text{many steps}$ &
  two-step correct proofs \\
\bottomrule
\end{tabular}
\end{center}
\medskip

In each case, the repair is not additional optimisation pressure
but evidence of admissible trajectories in the region the
training distribution had rendered sparse.

High correlation does not imply admissible substitution.
\end{proposition}

%=====================================================================
\section{Morphological Repair}
\label{sec:repair}
%=====================================================================

The easy-problem intervention is not a regulariser in the standard
sense of adding a constraint to the objective function. It is a
representational repair that restores admissible structure without
modifying the reward function.

This distinction matters because it points to a different theory of
what went wrong. A regulariser corrects an objective. A repair
corrects a representation. The baseline model's failure is not that
it is optimising the wrong objective---it is correctly maximising
expected reward under the verifiable binary reward function. The
failure is that its learned representation no longer faithfully
distinguishes between strategies whose downstream consequences differ.

The repair mechanism in admissibility theory takes the following
form. When a representation undergoes collapse by identifying two
distinct states $s_1$ and $s_2$ with the same representational
region, the minimal repair introduces a residual distinction that
separates the two states without abandoning their shared structure:

\[
\widetilde{F}(s_1) = (a, f), \qquad \widetilde{F}(s_2) = (b, f),
\qquad a \neq b.
\]

The shared component $f$ preserves whatever the two states genuinely
have in common (both are correct solutions). The residual components
$a \neq b$ restore the distinction that collapse had erased (one is
short, one is long; they are appropriate in different contexts).

Easy problems supply the $a \neq b$ signal. They are the training
examples for which $s_1$ (short correct solution) receives positive
reward and any attempt to produce $s_2$ (long solution) is either
truncated or wasteful. The group relative advantage assigns positive
weight to $s_1$ on easy problems in a way that the hard-only curriculum
never could, because hard problems genuinely require $s_2$. The repair
is morphological in the sense that it adds a minimal distinguishing
marker---easy-problem context---to the learned representation, not by
rewriting the underlying reward structure but by extending the evidence
base from which the policy infers which strategy to select.

\begin{theorem}[Easy Problems as Minimal Repair]
\label{thm:easyrepair}
Under hard-only training, the policy's learned representation satisfies
$F(s_1) \approx F(s_2)$ for short and long correct solutions
respectively. Reintroducing easy problems provides training examples
on which $s_1$ is rewarded and $s_2$ is penalised by truncation or
wasted budget. The GRPO advantage thereby assigns differential gradient
to $s_1$ and $s_2$ on easy problems, progressively separating the
two strategies in representation space. The repair is minimal: it does
not alter the representation of hard-problem solutions but restores
the short-solution fiber of the correctness projection.
\end{theorem}

\begin{proof}
By Theorem~\ref{thm:blind}, hard problems contribute zero differential
gradient between $s_1$ and $s_2$. By the budget-truncated reward
mechanism (formalised in Section~\ref{sec:advantage}), easy problems
produce positive reward for $s_1$ and zero or negative effective reward
for $s_2$ when $s_2$ exceeds the token budget or wastes allocation that
could serve harder problems. The GRPO advantage on easy problems is
therefore strictly positive for $s_1$ and non-positive for $s_2$.
Gradient descent under this signal separates $s_1$ from $s_2$ in the
policy's strategy representation. The separation is local to the
easy-problem regime: on hard problems the policy continues to produce
$s_2$ as before, since no new signal changes the hard-problem advantage.
The repair is therefore minimal in the sense that it modifies the
representation only in the short-solution fiber, not throughout
$\mathcal{S}$.
\end{proof}

A penalty-based approach to the same problem would compress the
long-solution region without expanding the short-solution region.
It makes $s_2$ costly without demonstrating that $s_1$ is admissible.
If the model has never produced short correct solutions under training,
penalising long solutions provides no gradient toward finding them.
The repair approach works by the opposite mechanism: it demonstrates
that $s_1$ is correct before the model has internalised that it cannot
be.

%=====================================================================
\section{The Geometry of Advantage Collapse}
\label{sec:advantage}
%=====================================================================

The GRPO advantage function is the formal site where projection
collapse manifests. We derive its behaviour under hard-only versus
mixed distributions and show that the hard-only curriculum is
structurally incapable of separating short from long correct
trajectories.

\subsection{Advantage blindness under hard-only training}

For a group of $G$ rollouts $\{y_i\}_{i=1}^G$ on query $x$, the
GRPO advantage is

\[
A_i = \frac{r(x, y_i) - \bar{r}}{\sigma_r},
\]

where $\bar{r}$ and $\sigma_r$ are the group mean and standard
deviation of rewards. Under binary rewards $r \in \{0,1\}$, if
$k$ rollouts are correct, $\bar{r} = k/G$ and
$\sigma_r = \sqrt{k(G-k)}/G$.

For a hard problem with per-rollout success probability $p \in (0,1)$:

\[
\mathbb{E}[A_i \mid r_i = 1] = \sqrt{\frac{1-p}{p}}, \qquad
\mathbb{E}[A_i \mid r_i = 0] = -\sqrt{\frac{p}{1-p}}.
\]

This advantage is assigned identically to all correct rollouts
regardless of length. For any two correct rollouts
$y_{\mathrm{short}}$ and $y_{\mathrm{long}}$:

\[
A_{\mathrm{short}} = A_{\mathrm{long}} = \sqrt{\frac{1-p}{p}}.
\]

\begin{theorem}[Advantage Blindness to Length]
\label{thm:blind}
Under the hard-only distribution, the GRPO advantage provides zero
differential signal between correct trajectories of different lengths.
For any $y_1, y_2$ with $r(x,y_1) = r(x,y_2) = 1$,

\[
A_1 - A_2 = 0
\]

regardless of $|\ell(y_1) - \ell(y_2)|$.
\end{theorem}

\begin{proof}
$A_i = (r_i - \bar{r})/\sigma_r$ depends only on $r_i$ and group
statistics $(\bar{r}, \sigma_r)$, which are functions of
$\{r_j\}_{j=1}^G$ alone. Since $r(x,y_1) = r(x,y_2) = 1$, both
rollouts contribute identically to all group statistics, so
$A_1 = A_2$.
\end{proof}

\begin{remark}
Theorem~\ref{thm:blind} is a structural result: the advantage
estimator cannot distinguish any property of a rollout beyond its
scalar reward. The gradient is therefore informationally incapable
of teaching the policy to prefer short correct solutions over long
ones when both exist. Repair must come from the training
distribution, not the objective function.
\end{remark}

\subsection{Advantage asymmetry under the mixed distribution}

Define the budget-truncated reward for budget $B$:

\[
r^B(x, y) = \begin{cases}
1 & \text{if } r(x,y) = 1 \text{ and } \ell(y) \leq B, \\
0 & \text{otherwise.}
\end{cases}
\]

\begin{theorem}[Effective Advantage Asymmetry Under Budget Constraints]
\label{thm:asymm}
Under the budget-truncated reward $r^B$, the combined effect of the
training distribution and the context limit produces an effective
advantage asymmetry on easy problems: when the fraction $q > 0$
of rollouts exceeds budget $B$, the effective learning signal
satisfies

\[
\mathbb{E}[A_{\mathrm{short}}] > \mathbb{E}[A_{\mathrm{long}}].
\]

This is a property of the training distribution and budget constraint
jointly, not of the GRPO objective in isolation.
\end{theorem}

\begin{proof}
Rollouts exceeding $B$ receive $r^B = 0$ by the truncation rule.
The group mean becomes $\bar{r} = k/G - q < k/G$ where $k$ counts
pre-truncation correct rollouts. A short correct rollout receives
reward~1, advantage $(1-\bar{r})/\sigma_{r^B} > 0$. A truncated
long rollout receives reward~0, advantage $-\bar{r}/\sigma_{r^B}
< 0$. Since $q > 0$ by hypothesis, the expected advantage
differential is strictly positive: $\mathbb{E}[A_{\mathrm{short}}]
- \mathbb{E}[A_{\mathrm{long}}] > 0$. The asymmetry arises from
the conjunction of (i) the fixed budget $B$ that converts long
rollouts into truncation events, and (ii) the easy-problem
distribution that ensures short correct rollouts exist within
$B$. Neither factor alone suffices.
\end{proof}

Together, Theorems~\ref{thm:blind} and~\ref{thm:asymm} characterise
the repair mechanism: hard problems provide zero differential gradient
between short and long correct solutions; easy problems provide
positive differential gradient favouring short solutions. The mixed
distribution has gradient structure the hard-only distribution
structurally lacks.

%=====================================================================
\section{Optimality of Repair over Penalty}
\label{sec:penalty}
%=====================================================================

\subsection{The penalty approach and its limit}

A length-penalised reward $r_\lambda(x,y) = r(x,y) - \lambda\ell(y)$
yields advantage differential

\[
A_{\mathrm{short}}^\lambda - A_{\mathrm{long}}^\lambda
= \frac{\lambda(\ell_{\mathrm{long}} - \ell_{\mathrm{short}})}{\sigma_{r_\lambda}} > 0.
\]

This creates gradient toward shorter solutions. It appears to solve
the problem. But it has an information-theoretic gap.

\begin{definition}[Admissibility Evidence]
Training example $(x,y)$ provides \emph{admissibility evidence}
for strategy $s$ if $r(x,y) = 1$ and $y$ instantiates $s$.
\end{definition}

\begin{theorem}[Information Gap of Length Penalties]
\label{thm:gap}
Under hard-only training with any $\lambda > 0$, admissibility
evidence for $s_{\mathrm{short}}$ is zero if the policy has not
yet discovered short correct solutions. Under the mixed distribution
with easy problems, admissibility evidence for $s_{\mathrm{short}}$
is positive regardless of the policy's prior behaviour.
\end{theorem}

\begin{proof}
Admissibility evidence for $s_{\mathrm{short}}$ requires $(x,y)$
with $r(x,y)=1$ and $\ell(y) \ll B$. Under hard-only training,
successful rollouts satisfy $\ell(y) \approx B$. If the policy has
not yet produced short correct solutions, no rollout in the hard-only
distribution satisfies both conditions. Adding $\lambda$ to the
reward does not change which rollouts the policy produces; it
only reweights rewards. Therefore admissibility evidence is zero
for any $\lambda$.

Under the mixed distribution, easy problems admit short correct
rollouts with $r = 1$ and $\ell \ll B$ with positive probability
even under a policy not yet specialised to brevity.
\end{proof}

\begin{corollary}[Repair Dominates Penalty in Early Training]
\label{cor:dominance}
Before the policy discovers short correct solutions, repair provides
strictly more information about $s_{\mathrm{short}}$ than any
penalty-augmented hard-only distribution.
\end{corollary}

Length penalties tell the policy that long solutions are costly.
Easy problems show the policy that short solutions are correct.
These are different pieces of information, and in the early training
phase only the second is available.

\begin{corollary}[Temporal Ordering]
\label{cor:temporal}
Repair is the appropriate early-phase intervention. Penalty
augmentation may consolidate an established brevity preference
in a later phase. Penalty before repair risks suppressing
exploration of $s_{\mathrm{short}}$ by making long solutions
costly without revealing that short ones exist.
\end{corollary}

%=====================================================================
\section{Fiber Support and Stability of the Repaired Representation}
\label{sec:topology}
%=====================================================================

\subsection{Formal fiber structure and the admissibility gap}

Let the policy strategy space decompose as
$\mathcal{S} = \mathcal{Q} \times \mathcal{L}$ with correctness
projection $\pi(q,\ell) = q$. The admissible fiber at $q^* =
\mathrm{correct}$ is $\mathcal{F}^* = \{q^*\} \times \mathcal{L}$.
The policy's empirical fiber $\hat{\mathcal{F}}$ contracts under
hard-only training:

\[
\hat{\mathcal{F}}_{\mathrm{hard}} \approx
\{q^*\} \times [\ell_{\min}^{\mathrm{hard}}, B],
\qquad \ell_{\min}^{\mathrm{hard}} \gg 0,
\]

and expands after repair:

\[
\hat{\mathcal{F}}_{\mathrm{repaired}} \approx
\{q^*\} \times [\ell_{\min}^{\mathrm{easy}}, B],
\qquad \ell_{\min}^{\mathrm{easy}} \ll \ell_{\min}^{\mathrm{hard}}.
\]

\begin{definition}[Fiber Admissibility Gap]
$\Delta\mathcal{F} = \mu(\mathcal{F}^* \setminus \hat{\mathcal{F}})$,
where $\mu$ weights lengths by their frequency in problem-appropriate
solutions. Admissibility failure is $\Delta\mathcal{F} > 0$.
\end{definition}

\subsection{Stability of the repaired fiber}

\begin{theorem}[Stability Condition]
\label{thm:stable}
Let $\rho_{\mathrm{easy}}$ be the fraction of easy problems in the
ongoing training distribution and $\delta > 0$ the threshold
gradient differential needed to sustain the short-solution region.
The repaired fiber is stable if and only if

\[
\rho_{\mathrm{easy}} \;\geq\; \rho^* :=
\frac{\delta}{\mathbb{E}_{x_e}\bigl[A_{\mathrm{short}}
- A_{\mathrm{long}} \mid x \text{ easy}\bigr]}.
\]
\end{theorem}

\begin{proof}
By Theorem~\ref{thm:blind}, hard problems contribute zero
differential gradient to the short-solution fiber. By
Theorem~\ref{thm:asymm}, each easy problem contributes expected
differential $\mathbb{E}[A_{\mathrm{short}} - A_{\mathrm{long}}]
> 0$. The net stabilising contribution is $\rho_{\mathrm{easy}}
\cdot \mathbb{E}[A_{\mathrm{short}} - A_{\mathrm{long}}]$.
Stability requires this to exceed $\delta$, giving the stated
bound.
\end{proof}

\begin{corollary}[A Testable Prediction]
\label{cor:prediction}
Setting $\rho_{\mathrm{easy}} = 0$ in Stage~2 (eliminating easy
problems entirely) violates the stability condition. The repair
framework predicts that extended Stage~2 training on hard-only
data will eventually regenerate verbosity, with re-collapse rate
determined by $\delta$ and the average advantage differential on
easy problems. This prediction is testable against extended
training runs.
\end{corollary}

\subsection{Monotone correspondence with reachability loss}

The fiber admissibility gap connects to the reachability loss
$\Lambda^\theta$ of the companion framework~\cite{flyxion2026}.
The merged strategy $s_{\mathrm{short}} \sqcup s_{\mathrm{long}}$
is the undifferentiated policy for ``correct solution'', and
reachability loss measures deployment contexts in which the merged
strategy is inferior to the differentiated one:

\[
\Lambda^\theta(s_{\mathrm{short}}, s_{\mathrm{long}})
= \mu^\theta\!\bigl(
\mathcal{A}_{s_{\mathrm{short}}}^\theta \cup
\mathcal{A}_{s_{\mathrm{long}}}^\theta
\bigr)
- \mu^\theta\!\bigl(
\mathcal{A}_{s_{\mathrm{short}} \sqcup s_{\mathrm{long}}}^\theta
\bigr).
\]

\begin{proposition}[Monotone Correspondence]
\label{prop:mono}
The fiber admissibility gap and reachability loss are monotonically
related: a larger gap cannot correspond to smaller loss. Formally,
for two policies with fiber gaps $\Delta\mathcal{F}_1 <
\Delta\mathcal{F}_2$,

\[
\Lambda^\theta_1 \;\leq\; \Lambda^\theta_2.
\]

In particular, $\Delta\mathcal{F} = 0$ implies $\Lambda^\theta = 0$.
\end{proposition}

\begin{proof}
$\Delta\mathcal{F}_2 > \Delta\mathcal{F}_1$ means the second policy
excludes a strictly larger region of $\mathcal{F}^*$ from its
empirical fiber. Futures in $\mathcal{F}^* \setminus
\hat{\mathcal{F}}_2$ but not in $\mathcal{F}^* \setminus
\hat{\mathcal{F}}_1$ are reachable under the first policy's
differentiated strategy but not under the second's. These futures
contribute to $\Lambda^\theta_2$ but not to $\Lambda^\theta_1$.
Therefore $\Lambda^\theta_2 \geq \Lambda^\theta_1$.

For the special case $\Delta\mathcal{F} = 0$: the empirical fiber
equals the admissible fiber, so the differentiated and merged
strategies access the same futures, giving $\Lambda^\theta = 0$.
\end{proof}

\begin{remark}
The monotone relation is sufficient for the essay's main claims
and easier to defend than exact equivalence. The two quantities
are defined over different formal objects---fiber gaps live in
strategy-length product space, reachability losses live in
admissible-futures measure space---and their relationship is
conceptual as well as formal. Monotonicity captures the
connection precisely without overclaiming identity.
\end{remark}

Stage~1 achieves $\Delta\mathcal{F} \approx 0$ and
$\Lambda^\theta \approx 0$ within the easy-problem domain.
Stage~2 maintains this while extending the policy into harder
domains, provided Theorem~\ref{thm:stable}'s condition holds.

%=====================================================================
\section{Two-Stage Repair and Capability Expansion}
\label{sec:stages}
%=====================================================================

The two-stage curriculum structure of the paper instantiates a
general principle of admissibility-based learning: repair before
expansion.

Stage~1 (emergent brevity) performs the morphological repair described
above. Easy problems restore the short-solution fiber of the
correctness projection. By the end of Stage~1, the model has learned
that length and quality are orthogonal dimensions of a correct solution,
not a single dimension.

Stage~2 (curriculum RLVR) performs capability expansion. Training
moves progressively from moderately difficult to highly difficult
problems. Crucially, Stage~2 does not reintroduce the collapse that
Stage~1 repaired. The reason is structural: the model now has a
representation in which correctness and length are separable. When
it encounters a hard problem that genuinely requires a long solution,
it extends its output because the problem requires it, not because
length is constitutive of correctness. The policy can modulate solution
length by problem structure rather than by a learned association with
a spurious proxy.

This sequence---repair first, then expand---is the correct order of
operations for any learning system that has undergone representational
collapse. Expansion before repair reintroduces the collapse, because
the harder domain's genuinely long solutions will again dominate the
training signal and contract the short-solution fiber. The paper's
empirical observation that Stage~2 brevity is preserved---minimum
response length increases modestly but average length does not revert
to baseline---confirms that the repair has been durable.

The two stages also correspond to two distinct forms of admissibility
maintenance. Stage~1 restores local faithfulness within the current
domain. Stage~2 extends admissibility into a new domain without
sacrificing the local faithfulness that Stage~1 established. Together
they implement what the domain divergence account of technical
vocabulary emergence predicts: a representation that is stable in
one domain can be extended into a harder domain without collapse,
provided the extension begins from the edge of the already-stable
region rather than from scratch.

%=====================================================================
\section{Scientific Repair and the Broader Pattern}
\label{sec:broader}
%=====================================================================

The repair structure observed in frugal reasoning is not specific to
language models. It recurs across domains where a representation has
undergone collapse under pressure from a biased sample of evidence.

A scientific anomaly repairs an ontology by forcing a distinction
between concepts previously treated as identical. Newtonian mechanics
treated mass and weight as proportional under terrestrial conditions,
which they are; the anomaly of inertial versus gravitational mass
forced the distinction open and revealed it to be theoretically
fundamental. The anomaly did not add new information to the domain.
It repopulated a previously empty fiber.

A biological mutation repairs an ecological niche by reopening a
previously inaccessible trajectory. A lineage that has become
specialised for one food source has undergone a representational
collapse of sorts---its phenotype no longer distinguishes between
environments where the specialised food is available and environments
where it is not. A mutation that restores dietary flexibility
expands the fiber of viable environments, restoring admissibility
without changing the organism's fundamental architecture.

A mathematical counterexample repairs a theorem by revealing a
hidden distinction. Cauchy's original proof of a theorem about
continuous functions treated pointwise and uniform continuity as
equivalent for the domain of interest; the counterexample forced
the distinction open and the theorem was restated correctly. The
counterexample did not contradict the original theorem under its
intended interpretation. It demonstrated that the intended
interpretation had collapsed a distinction the theorem required.

In every case, repair consists not of adding information but of
restoring admissible structure by demonstrating that a previously
empty fiber is non-empty. The easy mathematical problems are the
counterexamples that demonstrate the non-emptiness of the short
correct solution fiber. The repair is not a constraint. It is a
demonstration.

%=====================================================================
\section{The Broader Lesson for Machine Learning}
\label{sec:lesson}
%=====================================================================

The frugal reasoning result suggests a reframing of a class of
machine learning failures that are conventionally attributed to
insufficient data, insufficient model capacity, or insufficient
optimisation pressure.

Many such failures may instead arise from projection collapses
induced by unrepresentative training distributions. The model is
not failing to learn from the data it has. It is succeeding in
learning from the data it has, and the data happens to collapse
a distinction the task requires.

In such cases, the correct intervention is not more data of the
kind already present, not a stronger optimiser, and not a modified
objective function. The correct intervention is data that repopulates
the collapsed fiber---examples that demonstrate the admissibility of
the strategy the model has learned to exclude.

This principle applies beyond reasoning length. A model trained
on authoritative prose will conflate confidence with correctness.
A model trained on long arguments will conflate thoroughness with
validity. A model trained on formal language will conflate register
with precision. In each case, the model has undergone a projection
collapse: an accidental property of the training distribution has
been identified with an essential property of the task.

The repair in each case follows the same structure as the frugal
reasoning intervention. Introduce examples that demonstrate the
admissibility of the strategy the model has collapsed: confident
prose that is wrong, short arguments that are valid, informal
language that is precise. The collapse is repaired not by modifying
the objective but by expanding the fiber.

The lesson of \textit{Shorter but not Worse} is therefore not that
reasoning should be shorter. The lesson is that efficient reasoning
emerges when a system learns which distinctions matter.

Length is not reasoning.

Length is merely one trajectory through the space of reasoning.

A model becomes capable not when it learns to think longer, but when
it learns that thinking longer and thinking better are not the same
thing---and that the difference between them is exactly the
distinction that a well-curated training distribution must preserve.

%=====================================================================
\section{Conclusion}
\label{sec:conclusion}
%=====================================================================

Verbosity in reasoning models is a projection collapse. The hard-only
training curriculum distorts the fiber of the correctness projection,
contracting it from the full range of correct-solution lengths to a
region concentrated on long trajectories. The model does not fail to
optimise its objective; it succeeds too well, learning a representation
in which the accidental property of training-distribution solutions---
their length---has become identified with the essential property of
correct reasoning.

The information-theoretic account of why verbosity emerges is correct:
longer prefixes monotonically reduce conditional entropy of the final
answer, and a length-unpenalised policy has a weak incentive to exploit
this. But entropy reduction is not the goal. The goal is admissibility:
producing representations that keep the right futures reachable. A
model that reduces entropy by making all successful solutions look long
has made itself more predictable while simultaneously reducing its
navigational capacity.

The easy-problem repair restores admissibility by repopulating the
short-solution fiber of the correctness projection. It is not a
regulariser but a demonstration: it shows the model that short
solutions are admissible, that they lie within the fiber of correct
reasoning rather than outside it. The two-stage curriculum then
extends this repaired representation into harder domains without
reintroducing the collapse, because the policy has learned to
modulate solution length by problem structure rather than by a
spurious association with quality.

The pattern---collapse, repair, extension---is not specific to
language models. It recurs wherever a biased sample of evidence
causes a representation to treat an accidental property as
constitutive. Scientific anomalies, biological mutations, and
mathematical counterexamples all function as repair operators
in this sense: they demonstrate the non-emptiness of a fiber
that an accumulated evidence base had incorrectly treated as empty.

The measure of an intelligent system is not its ability to reduce
uncertainty. It is its ability to preserve the distinctions that
matter.

%=====================================================================
\begin{thebibliography}{99}

\bibitem{bounhar2026}
Bounhar, A., Abdine, H., Dufraisse, E., Chamma, A., Mohamed, A.,
Bouch, D., Vazirgiannis, M., \& Shang, G.\ (2026).
Shorter but not worse: Frugal reasoning via easy samples as length
regularizers in math RLVR.
\textit{arXiv:2511.01937v2}.

\bibitem{flyxion2026}
Flyxion (2026).
Why distinctions survive: A reachability theory of lexical preservation.
\textit{Independent Research Monograph}.

\bibitem{flyxion2026b}
Flyxion (2026).
The persistence hierarchy: Adjective order as reachability gradient.
\textit{Independent Research Monograph}.

\bibitem{flyxion2026latent}
Flyxion (2026).
Against latent fundamentalism: Admissibility and the limits of
dimensionality reduction.
\textit{Independent Research Monograph}.

\bibitem{shao2024}
Shao, Z., Wang, P., Zhu, Q., Xu, R., Song, J., Bi, X., Zhang, H.,
Zhang, M., Li, Y.\ K., Wu, Y., \& Guo, D.\ (2024).
DeepSeekMath: Pushing the limits of mathematical reasoning in open
language models.
\textit{arXiv:2402.03300}.

\end{thebibliography}

\end{document}
